
\content{
        \begin{definition} (Fair Division Method) A fair division method is a
        procedure that can be followed that will result in a division of the items
        in a way so that each party feels they have received their fair share.
        We make the following assumptions: 
        \begin{itemize}
        \item The parties are non-cooperative, so the method must operate
        without communication between parties. 
        \item The parties have no knowledge of what the other players like
        \item The parties act rationally, and do not make emotional decisions.
        \item The method should allow the parties to make a fair division
        without an outside arbitrator.
        \end{itemize}
        \end{definition}

        \begin{definition} (Fair Share) When $N$ parties divide something
        equally, each party's fair share is the amount they are entitled to.
        \end{definition}
}{% presentation
}{% nopresentation
}{% comments
Fair share is $1/N$.
}


\content{
        \begin{example} Suppose that 4 classmates are splitting equally a \$12
        pizza that is half pepperoni, half veggie that someone else bought them.
        What is each person's fair share? 
        \end{example}
}{% presentation
\vspace{1in}
}{% nopresentation
\vspace{1in}
}{% comments
Note that the dollar amount for each player is \$3, but each party might value
portions of the pizza differently, for ex: vegetarian...
}

\content{
        \begin{example} Suppose that in the above example, Steve likes pepperoni
        twice as much as veggie. Describe a fair share for Steve.
        \end{example}
}{% presentation
\vspace{2in}
}{% nopresentation
\vspace{2in}
}{% comments
Note that there are many different answers which are correct here.
}

\subsection{Divider Chooser Method}

\content{
        \begin{definition} (Divider-Chooser Method)
        \begin{enumerate}
        \item The divider cuts the item into two pieces that are, \textbf{in his
        eyes}, equal in value. 
        \item The chooser selects either of the two pieces. 
        \item the divider receives the remaining piece. 
        \end{enumerate}
        \end{definition}
}{% presentation
}{% nopresentation
\newpage
}{% comments
}

\content{
        \begin{example} Troy and Abed are sharing a \$12 Subway sandwich, which is
        a foot long and half veggie and half chicken. Troy likes chicken 2 times as much as
        veggies, and Abed likes chicken and veggies equally. If Troy is the
        divider, describe how he will split the sandwich (only vertical slices
        are allowed). Which part will Abed choose? Describe the perceived value
        of each piece to it's owner.
        \end{example}
}{%presentation
\vspace{4in}
}{% nopresentation
\vspace{4in}
}{% comments
Troy makes a cut 4.5'' into the chicken side. 
}

\content{
        \begin{example} Your turn! Annie and Jeff are now going to share a sandwich, which
        is half veggie and half chicken, and is a foot long. The sandwich costs
        \$16. Annie is a vegetarian, while 
        Jeff likes veggie 3 times as much as he likes chicken. If Jeff is the
        divider, describe how he will split the sandwich (only vertical slices
        are allowed). Which part will Annie choose? Describe the perceived value
        of each piece to it's owner.
        \end{example}
}{%presentation
\vspace{4in}
}{% nopresentation
\vspace{4in}
}{% comments
Jeff cuts 2/3 (4'') into the veggie side. 
}
